
\documentclass[template=tabling,81pt,headonall]{azmoon}
\usepackage{xepersian}
\usepackage{amsfonts}
\usepackage{graphicx}
\usepackage{svg}
\svgpath{ {./images/} }
\graphicspath{ {./images/} }
\settextfont{Yas}
\setdigitfont{A Iranian Sans}
\usepackage{fontawesome5}

\printanswers
    \teacher{محمد صالح علی اکبری}
    \teachertitle{دبیر}
    \city{گناباد}
    \schooltitle{دبستان}
    \school{پیشگامان ظهور}
    \grade{پنجم}
    \branch{۱}
    \topic{ریاضی}
    \examdate{مهر ۱۴۰۶}
    \answertime{۶۰ دقیقه}
    \begin{document}
	\begin{questions}
		\nointerlineskip%
		\vskip-\baselineskip
		\question[0.5]{%
در عدد هفت میلیارد و دویست و نود و شش میلیون و پانصد و چهل و سه هزار و پانصد و هفتاد، رقم مرتبه یکان هزار و مرتبه دهگان میلیون کدام گزینه است؟
    \begin{fourchoice}[2]\choice{یکان هزار ۳، دهگان میلیون ۹}
\choice{یکان هزار صفر، دهگان میلیون ۷}
\choice{یکان هزار ۳، دهگان میلیون ۵}
\choice{یکان هزار ۹، دهگان میلیون ۳}
\end{fourchoice}

            }\question[0.5]{%
کدام گزینه از خصوصیات لوزی نیست؟
    \begin{fourchoice}[1]\choice{مساحت لوزی مساوی حاصل‌ضرب دو قطر تقسیم بر ۲}
\choice{قطرها یکدیگر را نصف می‌کنند}
\choice{زاویه‌های روبه‌رو با هم مساوی هستند}
\choice{ضلع‌های روبه‌رو با هم مساوی نیستند}
\end{fourchoice}

            }\question[0.5]{%
گسترده عدد $3/61$ کدام گزینه است؟
    \begin{fourchoice}[2]\choice{$3 + 0/6 + 0/01$}
\choice{$3 + 6 + 0/01$}
\choice{$3 + 0/06 + 1$}
\choice{$36 + 0/1$}
\end{fourchoice}

            }\question[0.5]{%
اگر در یک کیسه سه تیله نارنجی و سه تیله سفید داشته باشیم، احتمال بیرون آمدن تیله سفید چقدر است؟
    \begin{fourchoice}[1]\choice{احتمال بیرون آمدن تیله سفید و نارنجی برابر است}
\choice{احتمال بیشتر تیله سفید در می‌آید}
\choice{احتمال کمتر تیله سفید در می‌آید}
\choice{احتمال خیلی زیاد تیله سفید در می‌آید}
\end{fourchoice}

            }\question[0.5]{%
کدام گزینه درست است؟
    \begin{fourchoice}[2]\choice{$0/25 = 0/250$}
\choice{$0/7 > 0/71$}
\choice{$0/306 > 0/36$}
\choice{$0/205 = 0/25$}
\end{fourchoice}

            }\section{درست یا نادرست بودن هر یک از عبارت‌های زیر را مشخص کنید.}\question[0.5]{%
مجموع زاویه‌های داخلی یک چهارضلعی $360$ درجه است. \hfill (\qquad\qquad)}\question[0.5]{%
برای به دست آوردن حجم یک مکعب مستطیل، طول و عرض و ارتفاع را با هم جمع می‌کنیم. \hfill (\qquad\qquad)}\question[0.5]{%
وقتی نیمساز یک زاویه $150$ درجه را رسم کنیم، دو زاویه $70$ درجه به دست می‌آید. \hfill (\qquad\qquad)}\question[0.5]{%
$3$ درصد شیر چربی است؛ یعنی از هر $100$ گرم شیر، $3$ گرم چربی است. \hfill (\qquad\qquad)}\question[0.5]{%
یک مربع فقط دو خط تقارن دارد. \hfill (\qquad\qquad)}\newpage\question[1]{%
اندازه زاویه (۱) و (۲) را پیدا کنید: \\ \\\\ \\\\ \\ زاویه $(1) = $ \\ زاویه $(2) = $‌
\\}\question[1]{%
در جاهای خالی عدد مناسب بنویسید.
    \begin{parts}[3]\part{$\dfrac{6}{5} = \dfrac{\framebox[0.7cm]{\vphantom{0}}}{25}$}
\part{$\dfrac{\framebox[0.7cm]{\vphantom{0}}}{21} = \dfrac{9}{7}$}
\part{$\dfrac{4}{3} = \dfrac{\framebox[0.7cm]{\vphantom{0}}}{12} = \dfrac{20}{\framebox[0.7cm]{\vphantom{0}}}$}
\end{parts}
‌
\\
    }\question[1]{%
پدر مریم $500$ تومان به او داد. مریم $60$ درصد پولش را یک جامدادی خرید. چند تومان برایش باقی مانده است؟‌
\\‌
\\‌
\\‌
\\}\question[1]{%
در این الگو شکل چهارم را رسم کنید و فرمول یا رابطه آن را بنویسید: \\ \\\\ \\\\ \\ رابطه یا فرمول = ‌
\\}\question[1.5]{%
جرم یک کیف مدرسه $1500$ گرم است. اگر در این کیف $3610$ گرم کتاب قرار دهیم، جرم کیف و کتاب‌ها با هم چقدر می‌شود؟‌
\\‌
\\‌
\\}\question[1.5]{%
محیط شکل زیر را پیدا کنید: \\ \\\\‌
\\‌
\\}\question[2]{%
دانش‌آموزان یک کلاس تصمیم گرفتند در جشن عید مبعث از دوستان خود با شربت پذیرایی کنند. اگر گنجایش یک لیوان $200$ سی‌سی باشد و تعداد دانش‌آموزان کلاس $25$ نفر باشد، برای تهیه شربت به یک ظرف چند لیتری نیاز است؟‌
\\‌
\\‌
\\}\question[2]{%
عملیات زیر را انجام دهید:
    \begin{LTR}
        \begin{parts}[1]\part{$4/26 \times 200 = $}
‌\\\part{$3 \div \dfrac{12}{8} = $}
‌\\\part{$3 \dfrac{2}{3} - 1 \dfrac{2}{5} = $}
‌\\\part{$7/1 - 3/025 = $}
\end{parts}
\end{LTR}
        ‌
\\
    }\end{questions}
    \end{document}
    